\chapter{Human Annotation Protocol}
\label{app:annotation}

This appendix records the realized human-reference protocol for the frozen PLT controlled benchmark. The purpose is to preserve the distinction between absolute cultural appropriateness and relative Best-of-4 preference.

\section*{Annotator-facing task}

For each evaluation item, the annotator was shown one prompt followed by four complete candidate responses labelled A, B, C and D. Generator identity, Vericult outputs, reward-model scores, direct-judge decisions, evidence traces and provisional labels were hidden.

For each candidate, the annotator recorded a three-level cultural-appropriateness rating corresponding operationally to inappropriate, partially appropriate and appropriate. After rating all four candidates, the annotator answered a forced comparative question selecting which of A, B, C or D was the most culturally appropriate response.

The forced winner is a relative judgment. It must not be interpreted as evidence that the chosen response is absolutely appropriate. Candidate-level ratings are therefore retained separately for the single-response evaluation.

\section*{Realized study}

The completed annotation contains five annotators, 30 prompts and four candidates per prompt. This yields 120 unique candidate responses, 600 candidate-level ratings and 150 Best-of-4 winner votes.

The realized form did \emph{not} include the earlier planned \emph{No clear winner} option or a 1--5 confidence field. Those fields are not reconstructed retrospectively. The frozen workbook is treated as immutable input; derived candidate, prompt and summary tables are regenerated deterministically from it and hashed in \texttt{data/evaluation/human\_gold\_manifest.json}.

\section*{Reference construction}

Candidate-level analysis preserves the full distribution of 1/2/3 ratings and reports mean, median and count summaries. Prompt-level relative analysis preserves the A--D vote distribution. A unique candidate with at least three of five votes is treated as a majority winner for resolved-winner comparison; prompts without such a majority remain ambiguous.

The frozen reference contains 17 prompts with a majority winner and 13 without one. Ten prompts tie for the largest winner-vote count; one prompt is unanimous. Fleiss' $\kappa$ over the complete winner matrix is 0.090494.

The aggregate is referred to throughout the thesis as a \emph{human reference judgment}, not cultural ground truth. A five-person panel cannot represent every culture or subgroup involved, and disagreement is itself informative for a culturally plural task.

\section*{Frozen annotation artifacts}

\begin{table}[htbp]
\centering
\footnotesize
\caption{Human-reference artifacts retained for reproducibility.}
\label{tab:annotation-artifacts}
\begin{tabularx}{\textwidth}{p{4.2cm}X}
\toprule
\textbf{Artifact} & \textbf{Role} \\
\midrule
\texttt{human\_annotations\_raw.xlsx} & Immutable survey export used as source data. \\
\texttt{human\_gold\_candidates.csv} & Candidate-level rating summaries for 120 responses. \\
\texttt{human\_gold\_prompts.csv} & Prompt-level A--D vote distributions and majority/ambiguity fields. \\
\texttt{human\_gold\_summary.json} & Aggregate counts and inter-rater statistics. \\
\texttt{human\_gold\_manifest.json} & SHA-256 hashes linking the raw workbook to all derived artifacts. \\
\bottomrule
\end{tabularx}
\end{table}

Human Gold is never loaded by Vericult or either baseline during inference. It is joined only after all machine predictions have been persisted.
